\chapter{Les stratégies des sociétés de gestion}\label{ch:s1:the-asset-managers-strategies}

Un indice est un portefeuille sur papier~: il se négocie sans frais, ne détient jamais de liquidités et se rééquilibre aux cours de clôture. Un fonds qui le réplique supporte chacun de ces
points, et Frino et Gallagher, qui ont étudié les \omterm{def:s1:index-rebalancing:fund}{fonds indiciels} adossés au S\&P 500, ont constaté que la tracking error est inévitable à cause des frictions de marché~; ils ont aussi constaté que les
\omterm{def:s1:index-rebalancing:fund}{fonds indiciels}, en moyenne, surpassent les fonds actifs après déduction des frais. Le métier des plus grandes sociétés de gestion consiste à maintenir un portefeuille proche d'un indice de référence à moindre coût,
à l'incliner délibérément, à transférer des capitaux d'un portefeuille à l'autre sans en perdre en chemin. Sur le marché synthétique de ce chapitre, détenir les 100 plus grandes
valeurs, parmi environ mille noms, avec des pondérations optimisées, suit l'indice à 1.64~\% par an, et les 400 plus grandes à 0.32~\%~; une transition portant sur 59~\% d'un fonds de 2~milliards de dollars
coûte 27.6~points de base si elle est réalisée en une journée, et au plus bas en coût plus risque si elle est étalée sur deux ou trois jours. La construction s'appuie sur \texttt{firm.assetmgr}.

\section{Indexation~: réplication et échantillonnage}

\begin{definition}[Réplication complète, réplication par échantillonnage]\label{def:s1:the-asset-managers-strategies:replication}
La \emph{réplication complète}\index{réplication complète} détient chaque composante d'un indice à son poids dans l'indice. \emph{Réplication par échantillonnage}%
\index{réplication par échantillonnage} détient un sous-ensemble des composantes, avec des pondérations choisies de sorte que les expositions aux risques du portefeuille soient celles de l'indice et que la tracking
error prévue soit aussi faible que possible, en échangeant des coûts moindres contre une part de tracking error.
\end{definition}

Un indice pondéré par la capitalisation des valeurs cotées de \texttt{firm.synthmkt} évolue de lui-même avec les cours~: le répliquer exactement n'exige des transactions que lorsque des valeurs entrent dans l'indice ou
en sortent, soit un taux de rotation à sens unique de 6.4~\% par an ici. L'échantillonnage exige un modèle de risque pour décider quelle valeur en représente quelle autre. Le chapitre utilise la
vérité propre au marché synthétique, le cas le plus favorable~: le bêta de chaque valeur sur le facteur de marché (avec sa variance conditionnelle), les facteurs sectoriels à 8~\% par an, les trois
styles à leurs volatilités injectées, avec des expositions à la date, et la volatilité spécifique de chaque valeur (le livre~7, chapitre~24, estime un tel modèle à partir des
rendements). Le \Cref{lst:s1:the-asset-managers-strategies:sampled} retient les $n$ plus grandes valeurs et résout un programme quadratique pour les pondérations à l'achat seul
qui minimisent la tracking error ex ante~; le portefeuille est rééquilibré tous les 63 jours et dérive avec les rendements dans l'intervalle.

\begin{center}
\small
\begin{tabular}{@{}lccccc@{}}
\toprule
noms détenus & 25 & 50 & 100 & 200 & 400\\
\midrule
tracking error réalisée & 4.84~\% & 2.96~\% & 1.64~\% & 0.81~\% & 0.32~\%\\
tracking error ex ante (moyenne) & 4.70~\% & 2.81~\% & 1.52~\% & 0.71~\% & 0.26~\%\\
taux de rotation à sens unique par an & 1.34 & 1.13 & 0.92 & 0.58 & 0.31\\
coût annuel à 10 pb par unité échangée (pb) & 13.4 & 11.3 & 9.2 & 5.8 & 3.1\\
\bottomrule
\end{tabular}
\end{center}

La tracking error baisse rapidement avec le nombre de valeurs (\cref{fig:s1:the-asset-managers-strategies:sampling}), de 40~\% à 60~\% à chaque doublement, parce que
chaque valeur ajoutée est plus petite et que le risque spécifique laissé de côté diminue avec le poids laissé de côté dans l'indice. Le chiffre ex ante sous-estime un peu le chiffre réalisé
même avec le vrai modèle~: entre deux rééquilibrages, les pondérations dérivent et les cotations changent. Les coûts évoluent à l'inverse de ce qu'on attendrait~: les petits échantillons
échangent davantage, parce que leurs pondérations doivent être réoptimisées à mesure que l'indice bouge. Sur les dix ans, l'échantillon de 25 noms a par hasard surpassé l'indice de 3.0~\% par
an~; avec une tracking error de 4.84~\%, cela relève du bruit, et un échantillon qui ne détient que les plus grandes valeurs constitue un pari sur les grandes capitalisations.

\begin{omfigure}
\begin{tikzpicture}
\begin{axis}[width=10cm,height=5.2cm,xmode=log,ymode=log,xlabel={\small noms détenus},ylabel={\small tracking error par an (\%)},
  tick label style={font=\footnotesize},xmin=20,xmax=500,ymin=0.2,ymax=6,grid=major,grid style={gray!25},legend pos=north east,
  legend style={font=\scriptsize},xtick={25,50,100,200,400},xticklabels={25,50,100,200,400},ytick={0.25,0.5,1,2,4},
  yticklabels={0.25,0.5,1,2,4},table/col sep=comma]
\addplot[omThm,thick,mark=*] table[x=names,y=te_pct]{figdata/strategies-1/29-the-asset-managers-strategies/sampling.csv};
\addplot[omDef,thick,dashed,mark=square*] table[x=names,y=te_ex_pct]{figdata/strategies-1/29-the-asset-managers-strategies/sampling.csv};
\legend{réalisée,ex ante}
\end{axis}
\end{tikzpicture}

\omcaption{\omterm{def:s1:the-asset-managers-strategies:replication}{Réplication par échantillonnage} de l'indice synthétique pondéré par la capitalisation~: tracking error réalisée et ex ante en fonction du nombre de plus grandes valeurs détenues, avec
des pondérations à l'achat seul minimisant la tracking error sous le vrai modèle de risque du marché, rééquilibrées tous les 63 jours (échelles logarithmiques). Données~:
\texttt{s1\_assetmgr.replicate}.}\label{fig:s1:the-asset-managers-strategies:sampling}
\end{omfigure}

Les \omterm{def:s1:index-rebalancing:fund}{fonds indiciels} réels ajoutent ce que le modèle laisse de côté~: les dividendes à réinvestir, les flux de capitaux des investisseurs, les changements d'indice annoncés à l'avance (chapitre~10), et
le choix du moment où les échanger. Frino, Gallagher et Oetomo ont constaté que les fonds passifs tirent profit d'un rééquilibrage moins rigide, et que les \omterm{def:s1:index-rebalancing:fund}{fonds indiciels} «~enrichis~»
rééquilibrent plus tôt et avec plus de patience autour des révisions d'indice et paient moins cher pour cela, et que les fonds passifs, lorsqu'ils s'écartent de l'indice, tendent à
surpondérer les valeurs les plus grandes et les plus liquides.

\section{Produits factoriels}

\begin{definition}[Produit factoriel]\label{def:s1:the-asset-managers-strategies:factor}
Un \emph{produit factoriel}\index{produit factoriel} est un portefeuille fondé sur des règles, généralement à l'achat seul, qui incline un indice large vers un style (valeur, momentum,
qualité, faible volatilité, taille) selon une règle publiée, commercialisé sous forme de fonds ou d'indice à des frais compris entre ceux d'un \omterm{def:s1:index-rebalancing:fund}{fonds indiciel} et ceux d'un gérant actif.
\end{definition}

Un \omterm{def:s1:the-asset-managers-strategies:factor}{produit factoriel} vend une exposition, non un rendement. Le chapitre incline l'indice par $\exp(kz)$ sur le z-score à la date du \omterm{def:s1:value-quality-and-low-risk:value}{ratio valeur comptable sur cours} de chaque valeur~:

\begin{center}
\small
\begin{tabular}{@{}lccc@{}}
\toprule
inclinaison $k$ & 0.25 & 0.5 & 1.0\\
\midrule
exposition valeur (en unités de z-score) & 0.26 & 0.52 & 1.03\\
tracking error & 1.57~\% & 3.08~\% & 6.00~\%\\
rendement excédentaire annuel & 0.07~\% & 0.07~\% & 0.06~\%\\
\bottomrule
\end{tabular}
\end{center}

L'exposition est livrée comme spécifiée, avec une tracking error proportionnelle à l'inclinaison~; le rendement, lui, ne l'est pas. Le marché synthétique injecte une prime de valeur,
mais un \omterm{def:s1:value-quality-and-low-risk:value}{ratio valeur comptable sur cours} mesuré avec le cours du jour est à découvert sur le momentum injecté (le chapitre~6 aboutit au même résultat), et les deux effets se compensent~: le produit ne rapporte presque
rien sur dix ans. L'acheteur d'un \omterm{def:s1:the-asset-managers-strategies:factor}{produit factoriel} a besoin d'une opinion sur le facteur, non sur le produit.

\section{Gestion relative à un indice de référence}

Les gérants actifs mesurés par rapport à un indice de référence évoluent dans l'espace que ces outils définissent. Leur risque est la tracking error, leur talent se mesure au ratio d'information
(livre~1, chapitre~3). La réponse moyenne-variance à la question de la quantité de tracking error à prendre est simple~: un gérant de ratio d'information IR, pénalisant la
variance active par un coefficient $\lambda$, maximise $\text{IR}\,\omega - \lambda\omega^2$ et retient $\omega = \text{IR}/(2\lambda)$~; Grinold et Kahn, dans
\emph{Active Portfolio Management} développe l'ensemble du cadre. L'indice échantillonné est l'extrémité sans talent de cette ligne, le \omterm{def:s1:the-asset-managers-strategies:factor}{produit factoriel}
une inclinaison constante, et le gérant actif une inclinaison variable dans le temps. Les outils du chapitre s'appliquent aux trois~: le même modèle de risque fixe la tracking error, le même
optimiseur maintient les expositions dans les limites, et les mêmes coûts décident de la fréquence du rééquilibrage.

\section{Gestion de transition}

\begin{definition}[Gestion de transition]\label{def:s1:the-asset-managers-strategies:transition}
La \emph{gestion de transition}\index{gestion de transition} consiste à faire passer un portefeuille d'un ensemble de positions à un autre (un nouveau
gérant, un nouvel indice de référence, une nouvelle allocation) au plus faible implementation shortfall, en arbitrant l'impact de marché d'un échange rapide contre le risque de détenir l'ancien
portefeuille pendant la construction du nouveau.
\end{definition}

Un fonds de pension qui remplace son échantillon indiciel de 100 noms par le produit value (une inclinaison de 0.5) doit échanger 59~\% d'un fonds de 2~milliards de dollars, dans un seul sens. L'échanger
en une journée coûte la demi-fourchette (5 points de base, supposés) et un impact en racine carrée sur la volatilité quotidienne de chaque valeur (\texttt{firm.tcost}, avec $\eta = 0.5$,
supposé)~; l'étaler sur davantage de jours réduit l'impact mais laisse la part non exécutée exposée au risque actif entre les deux portefeuilles, soit 3.46~\% par
an (21.8~points de base par jour). Le \Cref{lst:s1:the-asset-managers-strategies:sampled} inclut le planificateur.

\begin{center}
\small
\begin{tabular}{@{}lccccc@{}}
\toprule
jours pour achever & 1 & 2 & 3 & 5 & 10\\
\midrule
coût attendu (pb du fonds) & 27.6 & 21.3 & 18.5 & 15.6 & 12.8\\
écart type du résultat (pb) & 21.8 & 24.4 & 27.2 & 32.3 & 42.8\\
coût plus un écart type (pb) & 49.4 & 45.6 & 45.6 & 48.0 & 55.6\\
\bottomrule
\end{tabular}
\end{center}

Le compromis (\cref{fig:s1:the-asset-managers-strategies:transition}) présente un minimum large entre deux et trois jours pour un planificateur qui accorde autant de poids à un écart type de risque
qu'à un point de base de coût~; un planificateur qui ne se soucie que du coût attendu choisirait dix jours, avec un écart type
de 43~points de base. Les transitions réelles réduisent les deux termes en croisant les ordres~: les anciens et les nouveaux portefeuilles partagent souvent des valeurs, et les échanges entre eux, ou contre les flux des autres clients,
n'ont pas besoin de toucher le marché (livre~7, chapitre~27).

\begin{omfigure}
\begin{tikzpicture}
\begin{axis}[width=10cm,height=5.2cm,xlabel={\small jours pour achever la transition},ylabel={\small points de base du fonds},
  tick label style={font=\footnotesize},xmin=0,xmax=11,ymin=0,ymax=60,grid=major,grid style={gray!25},legend columns=3,
  legend style={font=\scriptsize,at={(0.5,-0.3)},anchor=north,draw=none,fill=none,/tikz/every even column/.append style={column sep=8pt}},
  table/col sep=comma]
\addplot[omThm,thick,mark=*] table[x=days,y=cost_bp]{figdata/strategies-1/29-the-asset-managers-strategies/transition.csv};
\addplot[omDef,thick,mark=square*] table[x=days,y=risk_bp]{figdata/strategies-1/29-the-asset-managers-strategies/transition.csv};
\addplot[black!60,thick,dashed,mark=triangle*] table[x=days,y=total_bp]{figdata/strategies-1/29-the-asset-managers-strategies/transition.csv};
\legend{coût attendu,écart type,coût plus un écart type}
\end{axis}
\end{tikzpicture}

\omcaption{Transition de 59~\% d'un fonds synthétique de 2~milliards de dollars, passant d'un échantillon indiciel de 100 valeurs à une inclinaison vers la valeur~: coût attendu
(demi-fourchette plus impact en racine carrée, paramètres supposés) et écart type dû à la détention de la part non exécutée, selon le nombre de jours
pris. Données~: \texttt{s1\_assetmgr.transition\_plan}.}\label{fig:s1:the-asset-managers-strategies:transition}
\end{omfigure}

\section{Fiches de stratégie}

\begin{strategyfile}{Réplication indicielle par échantillonnage}\label{strat:s1:the-asset-managers-strategies:index}
\sfield{Qui paie, et pourquoi} Des investisseurs qui paient des frais modestes pour le rendement du marché, à faible coût et à faible tracking error.
\sfield{Instruments et plateformes} Composantes de l'indice~; contrats à terme pour l'équitisation de la trésorerie.
\sfield{Signal} Aucun~: l'indice et un modèle de risque.
\sfield{Dimensionnement et exécution} \omterm{def:s1:the-asset-managers-strategies:replication}{Réplication complète} là où les coûts le permettent, échantillonnage sinon~; transactions à la clôture lors des changements d'indice, ou menées avec patience autour de ceux-ci.
\sfield{Coûts} Rotation due aux changements d'indice et à la réoptimisation.
\sfield{Comment elle s'éteint} Elle ne s'éteint pas~; les frais convergent vers zéro.
\sfield{Horizon, capacité, infrastructure} Permanent~; capacité très grande~; données d'indice, modèle de risque, optimiseur.
\sfield{Backtest honnête} L'indice tel qu'il était, avec ses changements et leurs dates d'annonce~; dividendes et trésorerie.
\sfield{Sources} Frino et Gallagher (2001)~; ce chapitre~: tracking error de 1.64~\% avec 100 valeurs, 0.32~\% avec 400.
\end{strategyfile}

\begin{strategyfile}{Indexation améliorée et produits factoriels}\label{strat:s1:the-asset-managers-strategies:enhanced}
\sfield{Qui paie, et pourquoi} Des investisseurs qui recherchent une prime factorielle ou de petits rendements excédentaires à un tracking error connu.
\sfield{Instruments et plateformes} Composantes de l'indice.
\sfield{Signal} Une règle d'inclinaison publiée, ou des transactions menées avec patience autour des changements d'indice.
\sfield{Dimensionnement et exécution} Inclinaisons dimensionnées selon un budget de tracking error~; rééquilibrage à échéance fixe.
\sfield{Coûts} Rotation modérée.
\sfield{Comment elle s'éteint} Des facteurs qui cessent de rapporter~; encombrement (crowding).
\sfield{Horizon, capacité, infrastructure} Des années~; grande capacité.
\sfield{Backtest honnête} Expositions point-in-time~; interaction du facteur avec les autres (valeur contre momentum).
\sfield{Sources} Frino, Gallagher et Oetomo (2005)~; ce chapitre~: exposition délivrée, rendement proche de zéro.
\end{strategyfile}

\begin{strategyfile}{Gestion de transition}\label{strat:s1:the-asset-managers-strategies:transition}
\sfield{Qui paie, et pourquoi} Des détenteurs d'actifs qui changent de gérants ou d'allocations, et qui paient pour éviter l'implementation shortfall d'une transition mal menée.
\sfield{Instruments et plateformes} Les positions des deux portefeuilles~; réseaux de croisement~; contrats à terme pour l'exposition intermédiaire.
\sfield{Signal} La liste des transactions et le risque actif entre les deux portefeuilles.
\sfield{Dimensionnement et exécution} Un calendrier qui arbitre impact et risque~; croisement en premier.
\sfield{Coûts} Fourchette, impact et risque de la part non exécutée.
\sfield{Comment elle s'éteint} Elle ne s'éteint pas~; elle se joue sur la qualité d'exécution.
\sfield{Horizon, capacité, infrastructure} Des jours~; un modèle de coûts et un modèle de risque.
\sfield{Backtest honnête} Implementation shortfall par rapport aux cours au moment de la décision, et non à l'achèvement.
\sfield{Sources} Ce chapitre~: 27.6 points de base en une journée, coût plus risque minimal sur deux ou trois jours.
\end{strategyfile}

\section{Tutoriel~: un savoir-faire invisible}\label{tut:s1:the-asset-managers-strategies:tut}

\begin{tutorial}
\textbf{Objectif.} Répliquer l'indice synthétique par échantillonnage, mesurer le tracking error et les coûts selon le nombre de valeurs, construire un produit value et planifier une
transition. \textbf{État final~:} les tableaux et les deux figures.
\begin{tutsteps}
\item \textbf{Échantillonnage, dérive, inclinaisons et transitions}.
\omcode{code/firm/assetmgr/firm_assetmgr.py}{44}{80}{Réplication par échantillonnage, dérive, inclinaisons et planificateur de transition.}\label{lst:s1:the-asset-managers-strategies:sampled}
\item \textbf{Le backtest de réplication}~: optimiser tous les 63 jours, laisser dériver entre-temps.
\omcode{code/strategies-1/29-the-asset-managers-strategies/python/s1_assetmgr.py}{65}{80}{Réplication par échantillonnage sur dix ans.}\label{lst:s1:the-asset-managers-strategies:replicate}
\item \textbf{Exécuter} \texttt{replicate(n)} pour cinq tailles, \texttt{full\_turnover()}, \texttt{value\_product(k)} et \texttt{transition\_plan()}, et
\texttt{fig\_assetmgr.py}.
\end{tutsteps}
\textbf{Pour aller plus loin.} Estimer le modèle de risque à partir des rendements (livre~7, chapitre~24) au lieu d'utiliser la vérité et comparer les tracking errors~; sélectionner
les valeurs en stratifiant par secteur plutôt que par taille~; croiser les noms communs à la transition avant d'échanger.
\end{tutorial}

\section{Construction~: la boîte à outils de la société de gestion}\label{bld:s1:the-asset-managers-strategies:assetmgr}

\begin{build}
\bfield{Rôle} Pondérations de l'indice, \omterm{def:s1:the-asset-managers-strategies:replication}{réplication par échantillonnage}, dérive des pondérations détenues, inclinaisons factorielles et planificateur de transition.
\bfield{Interface} \texttt{cap\_weights(price, shares, listed)}, \texttt{covariance(B, F, D)}, \texttt{sampled(b, Sigma, n)},
\texttt{drift\_active(R, w0, b\_path, start, end)}, \texttt{tilt(b, z, k)}, \texttt{transition(value, adv, sigma, half\_spread, eta, active\_sd, days)}.
\bfield{Règles} Pondérations positives seulement, de somme égale à un~; l'indice calculé à partir du cours de clôture précédent~; coûts issus de firm.tcost.
\bfield{Tests d'acceptation} \texttt{code/firm/assetmgr/tests/}~: pondérations d'indice et covariance calculées à la main, échantillonnage où toutes les valeurs donnent un tracking error nul et
avec dix valeurs un tracking error positif, dérive, inclinaison neutre, et coût de transition qui baisse tandis que le risque monte avec les jours.
\bfield{Pour aller plus loin} Échantillonnage stratifié~; modèle de risque estimé~; croisement des ordres dans les transitions.
\end{build}

\begin{omsources}
\item A.~Frino et D.~R.~Gallagher, ``Tracking S\&P 500 index funds'', \emph{Journal of Portfolio Management} 28(1), 2001.
\item A.~Frino, D.~R.~Gallagher et T.~N.~Oetomo, ``The index tracking strategies of passive and enhanced index equity funds'', \emph{Australian Journal
of Management} 30(1), 2005.
\item R.~C.~Grinold et R.~N.~Kahn, \emph{Active Portfolio Management}, 2e éd., McGraw-Hill, 2000.
\end{omsources}

\section{Exercices}

\begin{exercise}[$\star$]\label{exo:s1:the-asset-managers-strategies:1}
Si la tracking error diminuait avec la racine carrée du nombre de titres, de quel facteur doubler le nombre de titres la réduirait-elle~? Que montre le tableau entre 25 et 50
titres, et entre 100 et 200~?
\end{exercise}

\begin{exercise}[$\star$]\label{exo:s1:the-asset-managers-strategies:2}
Détenir exactement l'indice entraîne une rotation de 6.4~\% par an. À 10 points de base par unité échangée, combien cela coûte-t-il~?
\end{exercise}

\begin{exercise}[$\star$]\label{exo:s1:the-asset-managers-strategies:3}
Pourquoi la tracking error d'un \omterm{def:s1:the-asset-managers-strategies:factor}{produit factoriel} croît-elle avec l'inclinaison alors que son rendement n'est pas forcément croissant~?
\end{exercise}

\begin{exercise}[$\star\star$]\label{exo:s1:the-asset-managers-strategies:4}
Le risque actif de la transition est de 3.46~\% par an. Quel est-il par jour, et pourquoi s'agit-il de l'écart type d'une transition d'un jour~?
\end{exercise}

\begin{exercise}[$\star\star$]\label{exo:s1:the-asset-managers-strategies:5}
Pourquoi la tracking error réalisée d'un échantillon est-elle supérieure à sa tracking error ex ante, même avec le vrai modèle de risque~?
\end{exercise}

\begin{exercise}[$\star\star$]\label{exo:s1:the-asset-managers-strategies:6}
Avec une prime de valeur de 3~\% par an par unité d'exposition, que rapporterait une inclinaison d'exposition 0.52 si la valeur n'était corrélée à rien d'autre~?
Pourquoi n'a-t-elle rapporté que 0.07~\%~?
\end{exercise}

\begin{exercise}[$\star\star\star$]\label{exo:s1:the-asset-managers-strategies:7}
\emph{Programmation.} Modifier l'objectif de la transition pour qu'il soit le coût plus deux écarts types. Combien de jours le planificateur choisit-il maintenant~?
\end{exercise}

\begin{exercise}[$\star\star\star$]\label{exo:s1:the-asset-managers-strategies:8}
\emph{Trouver la faille.} «~Notre échantillon de 50 valeurs a battu l'indice de 2~\% par an pendant dix ans, donc l'échantillonnage crée de la valeur.~»
\end{exercise}

\section{Problème~: Artisanat invisible}

\begin{problem}[{Problème du week-end --- un indice, une inclinaison et une transition}]\label{pb:s1:the-asset-managers-strategies:1}
Le marché synthétique du chapitre et les données publiques.

\medskip\noindent\textbf{Partie I --- Indexation.}
\begin{enumerate}
\item Définir la \omterm{def:s1:the-asset-managers-strategies:replication}{réplication complète} et la \omterm{def:s1:the-asset-managers-strategies:replication}{réplication par échantillonnage}.
\item Qu'ont trouvé Frino et Gallagher au sujet des \omterm{def:s1:index-rebalancing:fund}{fonds indiciels}~?
\item Pourquoi détenir exactement l'indice exige-t-il si peu de transactions~?
\item Décrire le modèle de risque du chapitre et l'optimiseur d'échantillonnage.
\end{enumerate}

\medskip\noindent\textbf{Partie II --- Échantillonnage.}
\begin{enumerate}[resume]
\item Donner les tracking errors réalisées et ex ante selon le nombre de titres.
\item Pourquoi les petits échantillons échangent-ils davantage~?
\item Que signifiait le rendement excédentaire de l'échantillon de 25 titres~?
\item Qu'ont trouvé Frino, Gallagher et Oetomo au sujet du rééquilibrage~?
\end{enumerate}

\medskip\noindent\textbf{Partie III --- Produits et indices de référence.}
\begin{enumerate}[resume]
\item Définir un \omterm{def:s1:the-asset-managers-strategies:factor}{produit factoriel} et donner son exposition, sa tracking error et son rendement selon l'inclinaison.
\item Pourquoi le produit value n'a-t-il presque rien rapporté~?
\item Comment le cadre de Grinold et Kahn répartit-il la tracking error~?
\item Définir la \omterm{def:s1:the-asset-managers-strategies:transition}{gestion de transition}.
\end{enumerate}

\medskip\noindent\textbf{Partie IV --- Le verdict.}
\begin{enumerate}[resume]
\item Énoncer le \emph{résultat nommé}~: la tracking error de la \omterm{def:s1:the-asset-managers-strategies:replication}{réplication par échantillonnage} en fonction du nombre de titres détenus, et le shortfall de la transition.
\item Qu'est-ce qui limite la vitesse de la transition, et qu'est-ce qui limite sa patience~?
\item Comment le croisement des ordres réduit-il le coût d'une transition~?
\item Que changerait un modèle de risque estimé~?
\item Comment backtester honnêtement un \omterm{def:s1:index-rebalancing:fund}{fonds indiciel}~?
\item Quelle fiche de stratégie dépend le plus du modèle de risque~?
\item Quel lien ce chapitre entretient-il avec le chapitre~10~?
\item En une phrase~: en quoi consiste l'artisanat du gestionnaire d'actifs~?
\end{enumerate}
\end{problem}

\section{Questions d'entretien}

\begin{interviewq}[$\star$ \iqroles{chercheur}]\label{iq:s1:the-asset-managers-strategies:1}
Qu'est-ce qui cause la tracking error d'un \omterm{def:s1:index-rebalancing:fund}{fonds indiciel}~?
\end{interviewq}

\begin{interviewq}[$\star\star$ \iqroles{chercheur}]\label{iq:s1:the-asset-managers-strategies:2}
Comment choisiriez-vous les titres à détenir dans une \omterm{def:s1:the-asset-managers-strategies:replication}{réplication par échantillonnage} d'un indice de 3\,000 titres~?
\end{interviewq}

\begin{interviewq}[$\star\star$ \iqroles{trader}]\label{iq:s1:the-asset-managers-strategies:3}
Vous devez transférer un portefeuille de 5~milliards de dollars vers les positions du nouveau gestionnaire. Décrivez le plan.
\end{interviewq}

\begin{interviewq}[$\star\star$ \iqroles{développeur}]\label{iq:s1:the-asset-managers-strategies:4}
Quelles données et quels contrôles un \omterm{def:s1:index-rebalancing:fund}{fonds indiciel} doit-il prévoir autour d'un rééquilibrage de l'indice~?
\end{interviewq}

\begin{interviewq}[$\star\star$ \iqroles{risque}]\label{iq:s1:the-asset-managers-strategies:5}
Comment suivriez-vous le risque d'un \omterm{def:s1:the-asset-managers-strategies:factor}{produit factoriel}, et qu'est-ce qui vous inquiéterait~?
\end{interviewq}

\begin{interviewq}[$\star\star\star$ \iqroles{chercheur}]\label{iq:s1:the-asset-managers-strategies:6}
Une liste d'ordres est exécutée uniformément sur $N$ jours avec un coût attendu $c_0 N^{-1/2}$ (impact en racine carrée) et un risque $\sigma_a \sqrt{\sum_{d=0}^{N-1}(1 -
d/N)^2}$. Donnez une approximation du risque pour $N$ grand et trouvez le $N$ qui minimise le coût plus $\lambda$ fois le risque.
\end{interviewq}
